\begin{proof}[Proof of \cref{obj:thm:uniform-frontier}]
Fix \(0<\beta\le1\). We use the Legendre identities of
\cref{obj:lem:classical-legendre-facts} and the Hermite generating function and orthogonality of
\cref{obj:lem:classical-hermite-facts}. All constants introduced below depend at most on \(\beta\), unless explicitly indexed by an additional fixed parameter.

\begin{enumerate}
\item \textbf{Root facts used in the comparisons.}
For \(n\ge2\) and \(0\le\sigma\le1\), use the notation of
\cref{obj:def:uniform-rate}:
\[
\nu=2\beta+1,\qquad h_0=n^{-1/\nu}.
\]
Define, on \(0<h\le1\),
\[
\gamma_{\mathrm{bal}}(h)
=\log n+\nu\log h-\Psi(h,\sigma).
\]
We have
\[
0<h_0\le1,\qquad
\log n+\nu\log h_0=0.
\]
By \cref{obj:lem:noise-cost-scaling}, the function
\(\gamma_{\mathrm{bal}}\) is continuous. For \(0<h<t\le1\),
\[
\gamma_{\mathrm{bal}}(t)-\gamma_{\mathrm{bal}}(h)
=\nu\log(t/h)+\Psi(h,\sigma)-\Psi(t,\sigma)>0.
\]
Its endpoint values satisfy
\[
\gamma_{\mathrm{bal}}(h_0)=-\Psi(h_0,\sigma)\le0,
\qquad
\gamma_{\mathrm{bal}}(1)=\log n>0.
\]
The intermediate value theorem therefore gives a unique zero
\(\xi_{\mathrm{root}}\in[h_0,1]\). Strict increase also gives
\[
\{h\in[h_0,1]:\gamma_{\mathrm{bal}}(h)\le0\}
=[h_0,\xi_{\mathrm{root}}].
\]
Consequently, the supremum of this sublevel set equals
\(\xi_{\mathrm{root}}\), and the resolution in
\cref{obj:def:uniform-rate} satisfies
\[
h_\star(n,\sigma)=\xi_{\mathrm{root}},
\qquad
\log n+\nu\log h_\star(n,\sigma)
=\Psi(h_\star(n,\sigma),\sigma).
\]
Exponentiating yields the balance identity
\[
\exp\{\Psi(h_\star(n,\sigma),\sigma)\}
=n\,h_\star(n,\sigma)^\nu.
\]

If \(\sigma\le h_0\), then \(\Psi(h_0,\sigma)=0\), so uniqueness gives
\[
h_\star(n,\sigma)=h_0,\qquad
r_\beta(n,\sigma)=n^{-\beta/\nu}.
\]
This includes \(\sigma=0\).

If \(\sigma>h_0\), then \(\Psi(h_0,\sigma)>0\). Indeed, when
\(\sigma^4\le h_0\),
\[
(\sigma/h_0)^{2/3}>1.
\]
When \(h_0<\sigma^4\), we have \(h_0<1\) and
\(\sqrt{h_0}<\sigma^2\), whence
\[
h_0^{-1/2}>1,\qquad
1+\log(\sigma^2/\sqrt{h_0})\ge1,
\]
and their product exceeds one. Thus
\[
\gamma_{\mathrm{bal}}(h_0)<0,
\qquad
\gamma_{\mathrm{bal}}(\sigma)
=\log n+\nu\log\sigma
>\log n+\nu\log h_0=0.
\]
Strict increase now gives
\[
h_0<h_\star(n,\sigma)<\sigma.
\]
% lean: CausalSmith.Stat.RdTruesideNoiseFrontier.unique_rate_root
% lean: CausalSmith.Stat.RdTruesideNoiseFrontier.rateRoot_exp_balance
% lean: CausalSmith.Stat.RdTruesideNoiseFrontier.rateResolution_eq_direct_of_le
% lean: CausalSmith.Stat.RdTruesideNoiseFrontier.noisy_rateResolution_interior

\item \textbf{Lower comparisons.}
By \cref{obj:thm:direct-reduction}, there is \(c_{\mathrm{dir}}>0\) such that
\[
R_\beta(n,\sigma)\ge c_{\mathrm{dir}}n^{-\beta/\nu},
\qquad
\mathcal L_\beta(n,\sigma)\ge c_{\mathrm{dir}}n^{-\beta/\nu}
\]
for every admissible \(n,\sigma\). These are the required lower comparisons when
\(\sigma\le h_0\).

Consider \(\sigma>h_0\). For each integer \(m\ge2\), the scales of
\cref{obj:def:noise-support} satisfy
\[
b_m(\sigma)=\min\{1,\sigma\sqrt m/8\},
\qquad
h_m(\sigma)=b_m(\sigma)/m^2,
\qquad
0<b_m(\sigma)\le1.
\]
The support scale is nondecreasing in \(m\). Since
\(m^2\le4(m-1)^2\) for \(m\ge2\), it follows that
\[
h_{m-1}(\sigma)
=\frac{b_{m-1}(\sigma)}{(m-1)^2}
\le\frac{b_m(\sigma)}{(m-1)^2}
\le4h_m(\sigma).
\]
Also \(h_m(\sigma)\le m^{-2}\). In particular, any integer
\[
m>\sqrt{\frac{4096}{h_\star(n,\sigma)}}+2
\]
satisfies \(m\ge2\) and \(h_m(\sigma)\le h_\star(n,\sigma)/4096\).
Choose the smallest integer \(m\ge2\) with this latter property. At index two,
\[
b_2(\sigma)=\frac{\sigma\sqrt2}{8},
\qquad
h_2(\sigma)=\frac{\sigma\sqrt2}{32}
>\frac{h_\star(n,\sigma)}{4096},
\]
because \(h_\star(n,\sigma)<\sigma\) and
\(1\le\sqrt2<2\). Hence \(m\ge3\), and minimality gives
\[
h_{m-1}(\sigma)>\frac{h_\star(n,\sigma)}{4096}.
\]
Combining this with the predecessor bound gives
\[
\frac{h_\star(n,\sigma)}{16384}
<h_m(\sigma)
\le\frac{h_\star(n,\sigma)}{4096}.
\]

For this chosen index, define
\[
h_{\mathrm{low}}=h_m(\sigma),\qquad
b_{\mathrm{low}}=b_m(\sigma),\qquad
F_{\mathrm{cost}}(h)=1+\Psi(h,\sigma)\quad(0<h\le1).
\]
Thus \(0<h_{\mathrm{low}}\le h_\star(n,\sigma)<\sigma\le1\).
The ratio inequality of \cref{obj:lem:noise-cost-scaling} gives
\[
\frac{F_{\mathrm{cost}}(h_{\mathrm{low}})}
     {F_{\mathrm{cost}}(h_\star(n,\sigma))}
\ge
\left(\frac{h_\star(n,\sigma)}{h_{\mathrm{low}}}\right)^{1/2}
\ge64.
\]
Using \(h_{\mathrm{low}}^\nu\le h_\star(n,\sigma)^\nu\), the balance identity, and
\(\Psi(h_\star(n,\sigma),\sigma)\ge0\), we obtain
\[
\begin{aligned}
n h_{\mathrm{low}}^\nu
 \exp\{-F_{\mathrm{cost}}(h_{\mathrm{low}})/32\}
&\le
n h_\star(n,\sigma)^\nu
 \exp\{-2F_{\mathrm{cost}}(h_\star(n,\sigma))\}\\
&=\exp\{-2-\Psi(h_\star(n,\sigma),\sigma)\}\\
&\le e^{-2}.
\end{aligned}
\]

Define the observed alternative laws and their information budget by
\[
\mathsf P_\pm=(P^\pm_{b_{\mathrm{low}},m})_{\mathrm{obs}},
\qquad
\Delta_{\mathrm{info}}
=\frac{32}{3}\kappa^2h_{\mathrm{low}}^\nu
 \exp\{-F_{\mathrm{cost}}(h_{\mathrm{low}})/32\},
\qquad
\kappa=\frac1{100}.
\]
By \cref{obj:lem:cancellation-information-envelope},
\[
\chi^2(\mathsf P_+,\mathsf P_-)\le\Delta_{\mathrm{info}}<\infty,
\]
and the preceding calculation gives
\[
0\le n\Delta_{\mathrm{info}}
\le\frac{32}{3}\kappa^2e^{-2}
\le\frac{32}{30000}
\le\frac1{100}.
\]
Finiteness of this divergence supplies absolute continuity and an integrable
squared likelihood-ratio deviation. To spell out the product comparison, put
\[
\chi_{\mathrm{obs}}=\chi^2(\mathsf P_+,\mathsf P_-),
\qquad
\chi_{\mathrm{sample}}
=\chi^2(\mathsf P_+^{\otimes n},\mathsf P_-^{\otimes n}).
\]
The product likelihood identity and \(1+\chi_{\mathrm{obs}}\le e^{\chi_{\mathrm{obs}}}\) give
\[
1+\chi_{\mathrm{sample}}
=(1+\chi_{\mathrm{obs}})^n
\le e^{n\chi_{\mathrm{obs}}}
\le e^{1/100}
\le\frac1{1-1/100}
=\frac{100}{99}.
\]
Therefore
\[
\chi_{\mathrm{sample}}\le\frac1{99}\le\frac4{25},
\qquad
\operatorname{TV}(\mathsf P_+^{\otimes n},\mathsf P_-^{\otimes n})
\le\frac12\sqrt{\chi_{\mathrm{sample}}}
\le\frac15.
\]
Here the total-variation inequality follows by applying Cauchy--Schwarz to
the likelihood-ratio deviation; the positive and negative parts of that
deviation have equal mass.

The support and order conditions required by
\cref{obj:lem:two-point-transfer} hold. That result, which includes independent
procedure randomization, yields
\[
R_\beta(n,\sigma)\ge\frac45\kappa h_{\mathrm{low}}^\beta,
\qquad
\mathcal L_\beta(n,\sigma)\ge\frac65\kappa h_{\mathrm{low}}^\beta.
\]
Consequently, with
\[
c_{\mathrm{noise}}=\frac45\kappa\,16384^{-\beta},
\qquad
c_{\mathrm{mat}}=\min\{c_{\mathrm{dir}},c_{\mathrm{noise}}\}>0,
\]
the direct and noisy cases together give
\[
c_{\mathrm{mat}}r_\beta(n,\sigma)\le R_\beta(n,\sigma),
\qquad
c_{\mathrm{mat}}r_\beta(n,\sigma)\le\mathcal L_\beta(n,\sigma)
\]
uniformly over \(n\ge2\) and \(0\le\sigma\le1\).
% lean: CausalSmith.Stat.RdTruesideNoiseFrontier.exists_noiseResolution_bracket
% lean: CausalSmith.Stat.RdTruesideNoiseFrontier.rounded_noise_information_factor
% lean: CausalSmith.Stat.RdTruesideNoiseFrontier.sampleLaw_tv_le_one_fifth_of_chiSqExt
% lean: CausalSmith.Stat.RdTruesideNoiseFrontier.frontierRate_lower_bounds

\item \textbf{Upper comparison for the public radius.}
Choose \(K_{\mathrm{bias}}>0\) and \(D_{\mathrm{env}}>0\) from
\cref{obj:lem:positive-endpoint,obj:lem:variance-cost-envelope}, so that
\[
M_{L,\beta}\le K_{\mathrm{bias}}L^{-2\beta},
\qquad
V_L(\sigma)\le D_{\mathrm{env}}L^2
 \exp\{D_{\mathrm{env}}[1+\Psi(L^{-2},\sigma)]\}
\]
for every \(L\ge1\). Define
\[
\begin{aligned}
A_{\mathrm{dil}}&=4(D_{\mathrm{env}}+1)^2,\\
Q_{\mathrm{var}}&=\sqrt{160D_{\mathrm{env}}e^{D_{\mathrm{env}}+1}},\\
C_{\mathrm{up}}&=\max\left\{
 \frac12A_{\mathrm{dil}}^\beta,\,
 12K_{\mathrm{bias}}A_{\mathrm{dil}}^\beta+28Q_{\mathrm{var}}
 \right\}.
\end{aligned}
\]
These constants are positive, and \(A_{\mathrm{dil}}\ge4\).

The radii in \cref{obj:synth_6} and the public minimization in
\cref{obj:def:public-selection} give
\[
\rho_0=\frac12,\qquad
\rho_L=12M_{L,\beta}+28\sqrt{\frac{40V_L(\sigma)}n}\quad(L\ge1),
\qquad
E_\beta(n,\sigma)\le\rho_L\quad(0\le L\le L_{\max}),
\]
where the finite cap is
\[
L_{\max}=\left\lceil n^{1/(4\beta+2)}\right\rceil.
\]

If \(1\le A_{\mathrm{dil}}h_\star(n,\sigma)\), then
\[
E_\beta(n,\sigma)\le\frac12
\le\frac12A_{\mathrm{dil}}^\beta h_\star(n,\sigma)^\beta
\le C_{\mathrm{up}}r_\beta(n,\sigma).
\]
Suppose instead that \(A_{\mathrm{dil}}h_\star(n,\sigma)<1\), and define
\[
x_{\mathrm{up}}
=(A_{\mathrm{dil}}h_\star(n,\sigma))^{-1/2},
\qquad
L_{\mathrm{up}}=\lceil x_{\mathrm{up}}\rceil,
\qquad
h_{\mathrm{up}}=L_{\mathrm{up}}^{-2}.
\]
Since \(x_{\mathrm{up}}>1\), ceiling bounds give
\[
1\le L_{\mathrm{up}},
\qquad
x_{\mathrm{up}}\le L_{\mathrm{up}}
<x_{\mathrm{up}}+1\le2x_{\mathrm{up}}.
\]
Taking inverse squares yields
\[
\frac{A_{\mathrm{dil}}h_\star(n,\sigma)}4
\le h_{\mathrm{up}}
\le A_{\mathrm{dil}}h_\star(n,\sigma)<1.
\]
Moreover,
\[
x_{\mathrm{up}}
\le h_\star(n,\sigma)^{-1/2}
\le h_0^{-1/2}
=n^{1/(4\beta+2)},
\]
so monotonicity of the ceiling gives \(L_{\mathrm{up}}\le L_{\max}\).

Define the dilation multiplier by
\[
\alpha_{\mathrm{dil}}=\frac{h_{\mathrm{up}}}{h_\star(n,\sigma)}.
\]
The width bounds imply
\[
(D_{\mathrm{env}}+1)^2\le\alpha_{\mathrm{dil}}\le A_{\mathrm{dil}},
\qquad
\alpha_{\mathrm{dil}}\ge1,
\qquad
\alpha_{\mathrm{dil}}h_\star(n,\sigma)=h_{\mathrm{up}}\le1.
\]
Thus \cref{obj:lem:noise-cost-scaling} applies and gives
\[
1+\Psi(h_{\mathrm{up}},\sigma)
\le\max\left\{
1,\,
\alpha_{\mathrm{dil}}^{-1/2}
[1+\Psi(h_\star(n,\sigma),\sigma)]
\right\}.
\]
Because \(D_{\mathrm{env}}/\sqrt{\alpha_{\mathrm{dil}}}\le1\), multiplication by
\(D_{\mathrm{env}}\) gives
\[
\begin{aligned}
D_{\mathrm{env}}[1+\Psi(h_{\mathrm{up}},\sigma)]
&\le
\max\left\{
D_{\mathrm{env}},\,
1+\Psi(h_\star(n,\sigma),\sigma)
\right\}\\
&\le D_{\mathrm{env}}+1+\Psi(h_\star(n,\sigma),\sigma).
\end{aligned}
\]
Hence
\[
V_{L_{\mathrm{up}}}(\sigma)
\le D_{\mathrm{env}}L_{\mathrm{up}}^2
 e^{D_{\mathrm{env}}+1}
 \exp\{\Psi(h_\star(n,\sigma),\sigma)\}.
\]
The rounding bound also gives
\[
L_{\mathrm{up}}^2
\le\frac4{A_{\mathrm{dil}}h_\star(n,\sigma)}
\le\frac4{h_\star(n,\sigma)}.
\]
Using the root balance and \(\nu=2\beta+1\), we obtain
\[
\frac{40V_{L_{\mathrm{up}}}(\sigma)}n
\le160D_{\mathrm{env}}e^{D_{\mathrm{env}}+1}
 h_\star(n,\sigma)^{2\beta},
\]
and therefore
\[
\sqrt{\frac{40V_{L_{\mathrm{up}}}(\sigma)}n}
\le Q_{\mathrm{var}}h_\star(n,\sigma)^\beta.
\]
The bias bound gives
\[
M_{L_{\mathrm{up}},\beta}
\le K_{\mathrm{bias}}h_{\mathrm{up}}^\beta
\le K_{\mathrm{bias}}A_{\mathrm{dil}}^\beta
 h_\star(n,\sigma)^\beta.
\]
Since \(L_{\mathrm{up}}\) is an eligible candidate,
\[
\begin{aligned}
E_\beta(n,\sigma)
&\le\rho_{L_{\mathrm{up}}}\\
&\le
(12K_{\mathrm{bias}}A_{\mathrm{dil}}^\beta+28Q_{\mathrm{var}})
 h_\star(n,\sigma)^\beta\\
&\le C_{\mathrm{up}}r_\beta(n,\sigma).
\end{aligned}
\]
Both cases establish this comparison over the full parameter range.
% lean: CausalSmith.Stat.RdTruesideNoiseFrontier.upperDegree_width_bounds
% lean: CausalSmith.Stat.RdTruesideNoiseFrontier.upperDegree_le_Lmax
% lean: CausalSmith.Stat.RdTruesideNoiseFrontier.dilation_cost_absorption
% lean: CausalSmith.Stat.RdTruesideNoiseFrontier.certificate_le_frontierRate

\item \textbf{Matched estimation and interval chains.}
By \cref{obj:thm:finite-certificate}, the selected procedures satisfy
\[
\begin{aligned}
\sup_{P\in\mathcal M_\beta(\sigma)}
\mathbb E_{P,n}
|\widetilde\theta_{\beta,n,\sigma}-\theta(P)|
&\le E_\beta(n,\sigma),\\
\sup_{P\in\mathcal M_\beta(\sigma)}
\mathbb E_{P,n}\operatorname{length}(\widetilde I_{\beta,n,\sigma})
&\le2E_\beta(n,\sigma),\\
\widetilde I_{\beta,n,\sigma}&\in\mathcal I_{\beta,n,\sigma}.
\end{aligned}
\]
The minimax infima are bounded above by the objectives of these eligible procedures.

For completeness, the mean bounds imply
\[
|\theta(P)|=|\mu_1(P,0)-\mu_0(P,0)|\le\frac12.
\]
Thus the constant estimator zero has worst expected absolute error at most
\(1/2\), and the constant interval \([-1,1]\) has coverage one and length two.
Both extended minimax objectives are consequently finite. Conversion to real
numbers preserves their comparisons with the selected procedures.

Put
\[
C_{\mathrm{mat}}=2C_{\mathrm{up}}>0.
\]
Combining the lower comparisons, the finite certificate, and the upper comparison gives
\[
\begin{aligned}
c_{\mathrm{mat}}r_\beta(n,\sigma)
&\le R_\beta(n,\sigma)
\le\sup_{P\in\mathcal M_\beta(\sigma)}
 \mathbb E_{P,n}|\widetilde\theta_{\beta,n,\sigma}-\theta(P)|\\
&\le E_\beta(n,\sigma)
\le C_{\mathrm{mat}}r_\beta(n,\sigma),\\
c_{\mathrm{mat}}r_\beta(n,\sigma)
&\le\mathcal L_\beta(n,\sigma)
\le\sup_{P\in\mathcal M_\beta(\sigma)}
 \mathbb E_{P,n}\operatorname{length}(\widetilde I_{\beta,n,\sigma})\\
&\le2E_\beta(n,\sigma)
\le C_{\mathrm{mat}}r_\beta(n,\sigma).
\end{aligned}
\]
The lower transfer ranges over the stated randomized decision classes, with
arbitrary procedure degree.
% lean: CausalSmith.Stat.RdTruesideNoiseFrontier.matchedFrontier_of_scalar_bounds
% lean: CausalSmith.Stat.RdTruesideNoiseFrontier.matched_frontier

\item \textbf{Finite-sample branches and interfaces.}
Define the branch comparison constants by
\[
\begin{gathered}
A_{\mathrm{int}}=(1+3\nu/2)^{3/2},\qquad
c_\tau=\frac1{2(1+2\nu)},\qquad C_\tau=e+6,\\
c_{\mathrm{pow}}=C_\tau^{-2\beta},\qquad
C_{\mathrm{pow}}=c_\tau^{-2\beta},\\
c_{\mathrm{br}}=\min\{1,c_\tau,c_{\mathrm{pow}}\},
\qquad
C_{\mathrm{br}}=\max\{A_{\mathrm{int}},C_\tau,C_{\mathrm{pow}}\}.
\end{gathered}
\]
All are positive.

The direct branch was established above. In the noisy region,
\(0<\sigma\le1\) and \(h_0<h_\star(n,\sigma)<\sigma\). Evaluation at the interface gives
\[
\Psi(\sigma^4,\sigma)=\sigma^{-2}-1,
\qquad
\gamma_{\mathrm{bal}}(\sigma^4)
=\log(n\sigma^{4\nu})-(\sigma^{-2}-1).
\]
This identity also holds at \(\sigma=1\), where both quantities on the right
of the noise-cost identity are zero. Strict increase on \((0,1]\) gives
\[
\sigma^4\le h_\star(n,\sigma)
\quad\Longleftrightarrow\quad
\log(n\sigma^{4\nu})\le\sigma^{-2}-1.
\]
The comparison is valid throughout the noisy region, including when
\(\sigma^4<h_0\).

On this intermediate branch, define
\[
z=\left(\frac{\sigma}{h_\star(n,\sigma)}\right)^{2/3},
\qquad
S_{\mathrm{int}}=1+\log(n\sigma^\nu).
\]
The inequalities \(h_\star(n,\sigma)<\sigma\) and
\(\sigma^4\le h_\star(n,\sigma)\) imply
\[
1<z\le\sigma^{-2}.
\]
The intermediate noise-cost formula and the root identity give
\[
1+\Psi(h_\star(n,\sigma),\sigma)=z,
\qquad
\log z=\frac23\{\log\sigma-\log h_\star(n,\sigma)\}.
\]
Substitution yields
\[
z+\frac{3\nu}{2}\log z=S_{\mathrm{int}},
\qquad
h_\star(n,\sigma)=\sigma z^{-3/2}.
\]
The left side of the scalar equation is strictly increasing for \(z>0\),
so the indicated root is unique. Since \(z>1\),
\[
0\le\log z\le z,
\qquad
z\le S_{\mathrm{int}}\le(1+3\nu/2)z.
\]
Taking inverse powers gives
\[
\frac{\sigma}{S_{\mathrm{int}}^{3/2}}
\le h_\star(n,\sigma)
\le A_{\mathrm{int}}\frac{\sigma}{S_{\mathrm{int}}^{3/2}}.
\]

If the branch-test inequality is strict in the other direction, then
\(h_\star(n,\sigma)<\sigma^4\). Define
\[
\tau=h_\star(n,\sigma)^{-1/2}.
\]
We have
\[
\tau>\sigma^{-2}\ge1,\qquad
h_\star(n,\sigma)=\tau^{-2},
\]
and the compact cost formula gives
\[
1+\Psi(h_\star(n,\sigma),\sigma)
=\tau\{1+\log(\sigma^2\tau)\}.
\]
Since \(\log\tau=-\tfrac12\log h_\star(n,\sigma)\), the root equation becomes
\[
2\nu\log\tau+\tau\{1+\log(\sigma^2\tau)\}
=1+\log n.
\]
For \(\sigma^{-2}<\tau_1<\tau_2\), both
\(\log\tau\) and \(1+\log(\sigma^2\tau)\) increase strictly, and the latter
factor is positive. In particular,
\[
\begin{aligned}
\tau_1\{1+\log(\sigma^2\tau_1)\}
&<\tau_2\{1+\log(\sigma^2\tau_1)\}\\
&<\tau_2\{1+\log(\sigma^2\tau_2)\}.
\end{aligned}
\]
Thus the compact equation has a unique root on the indicated domain.

To compare that root with an explicit expression, define
\[
\Lambda_n=1+\log n,\qquad
\Omega_{n,\sigma}=\log(e+\sigma^2\Lambda_n),\qquad
y_{\mathrm{cmp}}=\sigma^2\tau,\qquad
q_{\mathrm{cmp}}=\frac{\Lambda_n}{\tau},\qquad
J_{\mathrm{cmp}}=1+\log y_{\mathrm{cmp}}.
\]
Here \(\Lambda_n>1\), \(\Omega_{n,\sigma}>0\),
\(y_{\mathrm{cmp}}>1\), and \(J_{\mathrm{cmp}}>1\). Dividing the compact
equation by \(\tau\) gives
\[
q_{\mathrm{cmp}}
=2\nu\frac{\log\tau}{\tau}+J_{\mathrm{cmp}}.
\]
Using \(0\le\log\tau\le\tau\) and \(\nu\le3\), we obtain
\[
J_{\mathrm{cmp}}\le q_{\mathrm{cmp}}
\le J_{\mathrm{cmp}}+2\nu
\le J_{\mathrm{cmp}}+6,
\qquad
q_{\mathrm{cmp}}\le(1+2\nu)J_{\mathrm{cmp}}.
\]
Also \(\sigma^2\Lambda_n=y_{\mathrm{cmp}}q_{\mathrm{cmp}}\), so
\[
\Omega_{n,\sigma}\ge1,\qquad
\Omega_{n,\sigma}\ge\log y_{\mathrm{cmp}},
\qquad
J_{\mathrm{cmp}}\le2\Omega_{n,\sigma}.
\]
For the reverse comparison,
\[
e+\sigma^2\Lambda_n
=e+y_{\mathrm{cmp}}q_{\mathrm{cmp}}
\le y_{\mathrm{cmp}}(e+J_{\mathrm{cmp}}+6).
\]
The elementary bound \(\log t\le t-1\) for \(t>0\) therefore gives
\[
\begin{aligned}
\Omega_{n,\sigma}
&\le\log y_{\mathrm{cmp}}+\log(e+J_{\mathrm{cmp}}+6)\\
&\le(J_{\mathrm{cmp}}-1)+(e+J_{\mathrm{cmp}}+5)\\
&=2J_{\mathrm{cmp}}+e+4\\
&\le(e+6)J_{\mathrm{cmp}}.
\end{aligned}
\]
It follows that
\[
q_{\mathrm{cmp}}\le2(1+2\nu)\Omega_{n,\sigma},
\qquad
\Omega_{n,\sigma}\le(e+6)q_{\mathrm{cmp}}.
\]
All quantities divided below are positive, and these inequalities yield
\[
c_\tau\frac{\Lambda_n}{\Omega_{n,\sigma}}
\le\tau
\le C_\tau\frac{\Lambda_n}{\Omega_{n,\sigma}}.
\]
Raising this comparison to the negative power \(-2\beta\) reverses the
inequalities. Since
\[
r_\beta(n,\sigma)=\tau^{-2\beta},
\qquad
\left(\frac{\Lambda_n}{\Omega_{n,\sigma}}\right)^{-2\beta}
=\left(\frac{\Omega_{n,\sigma}}{\Lambda_n}\right)^{2\beta},
\]
we obtain
\[
c_{\mathrm{pow}}
\left(\frac{\Omega_{n,\sigma}}{\Lambda_n}\right)^{2\beta}
\le r_\beta(n,\sigma)
\le C_{\mathrm{pow}}
\left(\frac{\Omega_{n,\sigma}}{\Lambda_n}\right)^{2\beta}.
\]
These are the asserted compact comparisons because
\(\Lambda_n=\log(en)\).

At the intermediate--compact equality,
\[
\gamma_{\mathrm{bal}}(\sigma^4)=0.
\]
If \(\sigma^4<h_0\), strict increase would give
\(0=\gamma_{\mathrm{bal}}(\sigma^4)<\gamma_{\mathrm{bal}}(h_0)\le0\).
Hence \(\sigma^4\in[h_0,1]\), and root uniqueness implies
\[
h_\star(n,\sigma)=\sigma^4.
\]
The equality condition can be written as
\[
\log n+4\nu\log\sigma=\sigma^{-2}-1.
\]
Using \(\log(\sigma^{-2})=-2\log\sigma\) and
\(\sigma^2\sigma^{-2}=1\), it gives both scalar identities
\[
\begin{aligned}
\sigma^{-2}+\frac{3\nu}{2}\log(\sigma^{-2})
&=1+\log(n\sigma^\nu),\\
2\nu\log(\sigma^{-2})
+\sigma^{-2}\{1+\log(\sigma^2\sigma^{-2})\}
&=1+\log n.
\end{aligned}
\]
Thus \(z=\tau=\sigma^{-2}\), and both resolution formulas equal \(\sigma^4\).

At the direct--noisy equality \(\sigma=h_0\), the direct result gives
\[
h_\star(n,\sigma)=\sigma=h_0,\qquad
\log(n\sigma^\nu)=0.
\]
The intermediate scalar equation then holds at \(z=1\), and
\(\sigma z^{-3/2}=\sigma\).

Finally, \(c_{\mathrm{br}}\le1,c_\tau,c_{\mathrm{pow}}\) and
\(C_{\mathrm{br}}\ge A_{\mathrm{int}},C_\tau,C_{\mathrm{pow}}\).
All comparison factors above are positive. Replacing their lower constants
by \(c_{\mathrm{br}}\) and their upper constants by \(C_{\mathrm{br}}\)
therefore gives simultaneous branch comparisons, with the exact equations
and interfaces unchanged.
% lean: CausalSmith.Stat.RdTruesideNoiseFrontier.intermediateCondition_iff
% lean: CausalSmith.Stat.RdTruesideNoiseFrontier.intermediate_branch_representation
% lean: CausalSmith.Stat.RdTruesideNoiseFrontier.intermediate_branch_bounds
% lean: CausalSmith.Stat.RdTruesideNoiseFrontier.compact_branch_representation
% lean: CausalSmith.Stat.RdTruesideNoiseFrontier.compactScalar_bounds
% lean: CausalSmith.Stat.RdTruesideNoiseFrontier.compactPower_bounds
% lean: CausalSmith.Stat.RdTruesideNoiseFrontier.intermediate_compact_interface
% lean: CausalSmith.Stat.RdTruesideNoiseFrontier.direct_interface
% lean: CausalSmith.Stat.RdTruesideNoiseFrontier.frontier_branches

\item \textbf{Shrinking-noise sequences.}
Fix \(0<a<1/\nu\) and set \(\sigma=n^{-a}\). Define
\[
\eta_a=1-4a\nu,\qquad
\epsilon_a=\frac1{2(|\eta_a|+1)}.
\]
Since \(2a>0\), the relation \(\log n=o(n^{2a})\) and divergence of \(n^{2a}\)
give an integer \(N\ge2\) such that, for all \(n\ge N\),
\[
\log n\le\epsilon_a n^{2a},
\qquad n^{2a}\ge2.
\]
Consequently,
\[
\eta_a\log n
\le|\eta_a|\epsilon_a n^{2a}
\le\frac12n^{2a}
\le n^{2a}-1.
\]
The relevant quantities are exactly
\[
\log(n\sigma^{4\nu})=(1-4a\nu)\log n,
\qquad
\sigma^{-2}=n^{2a}.
\]
Thus the intermediate branch holds for \(n\ge N\); moreover,
\(0<\sigma\le1\) and \(\sigma>h_0\) for every \(n\ge2\).

Define
\[
d_{\mathrm{shrink}}=1-a\nu>0,\qquad
K_{\mathrm{shrink}}=d_{\mathrm{shrink}}+\frac1{\log2}>0.
\]
For \(n\ge2\),
\[
1+\log(n\sigma^\nu)=1+d_{\mathrm{shrink}}\log n,
\]
and \(\log n\ge\log2\) implies
\[
d_{\mathrm{shrink}}\log n
\le1+\log(n\sigma^\nu)
\le K_{\mathrm{shrink}}\log n.
\]
Applying the intermediate comparison with \(c_{\mathrm{br}},C_{\mathrm{br}}\)
therefore gives, for \(n\ge N\),
\[
\frac{c_{\mathrm{br}}}{K_{\mathrm{shrink}}^{3/2}}
\,n^{-a}(\log n)^{-3/2}
\le h_\star(n,n^{-a})
\le
\frac{C_{\mathrm{br}}}{d_{\mathrm{shrink}}^{3/2}}
\,n^{-a}(\log n)^{-3/2}.
\]
Taking the positive power \(\beta\) yields
\[
\begin{aligned}
\left(\frac{c_{\mathrm{br}}}{K_{\mathrm{shrink}}^{3/2}}\right)^\beta
n^{-a\beta}(\log n)^{-3\beta/2}
&\le r_\beta(n,n^{-a})\\
&\le
\left(\frac{C_{\mathrm{br}}}{d_{\mathrm{shrink}}^{3/2}}\right)^\beta
n^{-a\beta}(\log n)^{-3\beta/2}.
\end{aligned}
\]
Both constants are positive and may depend on \(a\).

For \(a\ge1/\nu\), monotonicity in the exponent for \(n>1\) gives
\[
0<n^{-a}\le n^{-1/\nu}=h_0\le1.
\]
Hence the exact direct branch applies for every \(n\ge2\), with
\[
h_\star(n,n^{-a})=h_0,
\qquad
r_\beta(n,n^{-a})=n^{-\beta/\nu}.
\]
% lean: CausalSmith.Stat.RdTruesideNoiseFrontier.eventually_shrinking_intermediate
% lean: CausalSmith.Stat.RdTruesideNoiseFrontier.shrinking_intermediate_rate_bounds
% lean: CausalSmith.Stat.RdTruesideNoiseFrontier.frontier_sequences

\item \textbf{Fixed positive noise.}
Fix \(0<\sigma\le1\), and define
\[
\begin{aligned}
\ell_{\mathrm{noise}}&=\log(\sigma^2),\\
M_{\mathrm{noise}}&=\max\left\{
1,\,-2\ell_{\mathrm{noise}},\,|\log(e+\sigma^2)|
\right\},\\
K_{\mathrm{noise}}&=1+|\log(e+\sigma^2)|,\\
T_{\mathrm{noise}}&=\sigma^{-2}-1-4\nu\log\sigma.
\end{aligned}
\]
Since \(\log n\to\infty\) and \(h_0\to0\), there is \(N\ge2\) such that, for
\(n\ge N\),
\[
e^{M_{\mathrm{noise}}}\le\log n,\qquad
T_{\mathrm{noise}}+1\le\log n,\qquad
h_0<\sigma.
\]
The middle inequality gives
\[
\log(n\sigma^{4\nu})
=\log n+4\nu\log\sigma
>\sigma^{-2}-1,
\]
so the compact branch applies.

Retain
\[
\Lambda_n=1+\log n,\qquad
\Omega_{n,\sigma}=\log(e+\sigma^2\Lambda_n),
\]
and define
\[
L_{\mathrm{log}}=\log\Lambda_n.
\]
The first threshold implies
\[
L_{\mathrm{log}}\ge M_{\mathrm{noise}}\ge1,
\qquad
L_{\mathrm{log}}\ge-2\ell_{\mathrm{noise}}.
\]
Monotonicity of the logarithm gives
\[
\Omega_{n,\sigma}
\ge\log(\sigma^2\Lambda_n)
=\ell_{\mathrm{noise}}+L_{\mathrm{log}}
\ge\frac12L_{\mathrm{log}}.
\]
Since \(\Lambda_n\ge1\),
\[
e+\sigma^2\Lambda_n\le(e+\sigma^2)\Lambda_n,
\]
and hence
\[
\begin{aligned}
\Omega_{n,\sigma}
&\le\log(e+\sigma^2)+L_{\mathrm{log}}\\
&\le|\log(e+\sigma^2)|L_{\mathrm{log}}+L_{\mathrm{log}}\\
&=K_{\mathrm{noise}}L_{\mathrm{log}}.
\end{aligned}
\]
Applying the compact rate comparison and taking the power \(2\beta>0\) gives
\[
\begin{aligned}
c_{\mathrm{br}}2^{-2\beta}
\left(\frac{L_{\mathrm{log}}}{\Lambda_n}\right)^{2\beta}
&\le r_\beta(n,\sigma)\\
&\le
C_{\mathrm{br}}K_{\mathrm{noise}}^{2\beta}
\left(\frac{L_{\mathrm{log}}}{\Lambda_n}\right)^{2\beta}.
\end{aligned}
\]
Since \(\Lambda_n=\log(en)\) and \(L_{\mathrm{log}}=\log\log(en)\), this is the
fixed-noise assertion with the permitted dependence on \(\sigma\).
% lean: CausalSmith.Stat.RdTruesideNoiseFrontier.fixed_noise_rate_bounds
% lean: CausalSmith.Stat.RdTruesideNoiseFrontier.frontier_sequences

\item \textbf{The case \(\beta=1\), \(\sigma=n^{-1/6}\).}
For every \(n\ge2\),
\[
\nu=3,\qquad h_0=n^{-1/3},\qquad
\sigma=n^{-1/6}>h_0,\qquad 0<\sigma\le1.
\]
Power identities give
\[
\sigma^{12}=n^{-2},\qquad
n\sigma^{12}=n^{-1},\qquad
\sigma^{-2}=n^{1/3}.
\]
Thus
\[
\log(n\sigma^{12})=-\log n
\le n^{1/3}-1=\sigma^{-2}-1,
\]
so the intermediate branch holds at every such sample size. Also
\[
n\sigma^3=n^{1/2},
\qquad
\log(n\sigma^3)=\frac12\log n,
\qquad
r_1(n,\sigma)=h_\star(n,\sigma).
\]
The intermediate comparison at \(\beta=1\) therefore yields
\[
r_1(n,n^{-1/6})
=h_\star(n,n^{-1/6})
\asymp
\frac{n^{-1/6}}{[1+\tfrac12\log n]^{3/2}},
\]
with constants independent of \(n\).
% lean: CausalSmith.Stat.RdTruesideNoiseFrontier.beta_one_frontier

\item \textbf{A simultaneous choice of constants.}
Finally, define
\[
c=\min\{c_{\mathrm{mat}},c_{\mathrm{br}}\},
\qquad
C=\max\{C_{\mathrm{mat}},C_{\mathrm{br}}\}.
\]
Then
\[
0<c,\qquad 0<C,\qquad
c\le c_{\mathrm{mat}},c_{\mathrm{br}},
\qquad
C\ge C_{\mathrm{mat}},C_{\mathrm{br}}.
\]
The root bounds give
\[
h_\star(n,\sigma)\ge h_0>0,
\qquad
r_\beta(n,\sigma)=h_\star(n,\sigma)^\beta\ge0.
\]
Thus, in the matched chains,
\[
c\,r_\beta(n,\sigma)
\le c_{\mathrm{mat}}r_\beta(n,\sigma),
\qquad
C_{\mathrm{mat}}r_\beta(n,\sigma)
\le C\,r_\beta(n,\sigma).
\]
The branch comparisons admit the same replacement: their factors
\[
\frac{\sigma}{[1+\log(n\sigma^\nu)]^{3/2}},
\qquad
\frac{\Lambda_n}{\Omega_{n,\sigma}},
\qquad
\left(\frac{\Omega_{n,\sigma}}{\Lambda_n}\right)^{2\beta}
\]
are positive on their respective noisy branches. Multiplication therefore
preserves the inequalities when their lower constants are reduced to \(c\)
and their upper constants are increased to \(C\). The root, branch equations,
uniqueness assertions, and interfaces retain their exact values. Together
with the sequence comparisons and the case \(\beta=1\), this proves every
assertion simultaneously.
% lean: CausalSmith.Stat.RdTruesideNoiseFrontier.matched_frontier_weaken
% lean: CausalSmith.Stat.RdTruesideNoiseFrontier.frontier_branches_weaken
% lean: CausalSmith.Stat.RdTruesideNoiseFrontier.uniform_frontier
\end{enumerate}
\end{proof}
